\begin{table}[htbp]
\centering\small
\begin{tabularx}{\linewidth}{rlXrrrrr}
\toprule
順 & 型 & テーマ & 問いの価値 & 信頼性 & 実行可能性 & 競争と出口 & 合計 \\
\midrule
1 & 手法 & T031 LLM が作った変数を回帰に使うときの測定誤差の補正（T033 を統合） & 65 & 80 & 100 & 65 & 76.2 \\
2 & 手法 & T057 AI の仕様探索は偽陽性をどれだけ増やすか：取り置きデータによる測定（T056 を統合） & 70 & 73 & 95 & 60 & 73.5 \\
4 & 構造 & T069 ニューラルネットによる構造パラメータの推定：精度・速度・推論の検証 & 60 & 73 & 90 & 60 & 69.8 \\
5 & 手法 & T035 統計書・行政文書の表を LLM で構造化するときの精度と誤差の伝播 & 55 & 67 & 85 & 75 & 69.5 \\
7 & 手法 & T016 回帰不連続の操作検定の検出力：どの程度の操作なら見逃すか & 55 & 73 & 95 & 55 & 67.8 \\
8 & 手法 & T004 処置の強さが連続のときの差の差：推定と診断の実装 & 60 & 67 & 95 & 55 & 67.6 \\
9 & 手法 & T025 実データに似せたシミュレーションで推定量を選ぶ手順の妥当性 & 60 & 67 & 95 & 55 & 67.6 \\
11 & 因果 & T044 クラスタの水準の選択が結論を変える頻度：公刊論文の再点検 & 60 & 73 & 85 & 55 & 67.3 \\
15 & 因果 & T006 ゼロを含む結果変数の対数変換が結論を変える頻度の大規模点検 & 55 & 80 & 85 & 50 & 65.8 \\
19 & 構造 & T068 観測されない異質性をもつ動学モデルの推定法の小標本での性能 & 55 & 67 & 90 & 55 & 65.3 \\
\bottomrule
\end{tabularx}
\caption{上位10個の案（AI の推奨）。順は100件の中での合計点の順位。データは \texttt{theme/research/recommended.csv} と \texttt{scored.csv}。}
\label{tab:th-top10}
\end{table}

\begin{table}[htbp]
\centering\small
\begin{tabularx}{\linewidth}{lXX}
\toprule
テーマ & 推す理由 & いちばんの弱点・未確認の点 \\
\midrule
T031 & LLM で作った変数を回帰に入れる研究が増えており、使う場面が広い。合成データだけで偏りの大きさと補正の効き目を示せる。4群のうち3群が上位 & 偏りの補正は Battaglia らが提案済みである。統計側の手法（設計にもとづく補正、予測を使った推論）との差を示せるかが勝負になる \\
T057 & このラボのロックと仕様の記録が、そのまま測定の道具になる。効果が無いと分かっている設定で偽陽性率を測れば、どちらの結果でも知見になる & 2026年に近い研究が出始めた。エージェントへの指示の書き方に結果が左右される \\
T069 & 構造推定のトラックで最上位。古典的なモデルで従来法と比べ、どこで失敗するかを示す検証は、提案者以外がやる価値がある & 手法の進歩が速い。提案者（Wei・Jiang）の検証が先に広がるおそれがある \\
T035 & 日本の古い統計書は手つかずで、作ったデータが次の研究の資産になる。「競争と出口」が75点で最も高い & 文書ごとに読み取りの精度が大きく違う。照合に使う正解データをどこまで用意できるかは未確認 \\
T016 & 検定を通ったことを安全の根拠にする慣行に、直接効く。合成データで完結し、最初の結果が数週間で出る & 問いが小さい。論文としては短報どまりになりうる（「問いの価値」55点） \\
T004 & 処置の強さが連続の場合は、使いやすい実装と診断がまだ薄い。ライブラリを出口にできる & 理論が発展途上で、推奨が変わりうる。提案者自身のパッケージ（contdid）が先行している \\
T025 & このラボの段階6（シミュレーター）そのものが主題になる。自分たちの手順の妥当性を確かめる研究でもある & Advani らが既に否定的な結果を出している。同じ結論の追認に終わるおそれがある \\
T044 & 因果推論のトラックで最上位。理論は示されたが、公刊論文で結論が変わる頻度を数えた研究は検索で見つからなかった。再推定の量で勝負できる & 適切なクラスタの水準を決めるには、論文ごとに制度の知識が要る \\
T006 & 因果推論のトラックの2番目。対数変換の問題は Chen・Roth が示しており、公刊論文での頻度を数える作業は AI の量が効く & 問題提起は済んでおり、追加の知見が薄いとブラインド評者が指摘した。合計点は T030（感度分析の一括適用）と同点である \\
T068 & 構造推定のトラックの2番目。実装が難しく参入が少ない。合成データで完結する & 需要が細い。比較研究の域を出ないおそれがある（「問いの価値」「競争と出口」が55点） \\
\bottomrule
\end{tabularx}
\caption{上位10個の案の根拠と弱点。近い既存研究・データ・競合の出典は \texttt{theme/research/themes.csv} の各欄にある。}
\label{tab:th-top10why}
\end{table}
